\section{Objetivo e hipótesis}

\subsection{Objetivo general}

Comparar los mecanismos de comunicación entre procesos sockets TCP y gRPC
mediante su implementación en un sistema distribuido de tres nodos, la ejecución
de pruebas repetidas de latencia y rendimiento, y el registro de eventos con
relojes lógicos de Lamport, para evaluar su comportamiento bajo distintas cargas
y condiciones de concurrencia, y verificar el orden causal de los mensajes.

\subsection{Objetivos específicos}

\begin{enumerate}
    \item Implementar ambos mecanismos con mensajes equivalentes y un reloj
    lógico de Lamport por nodo.
    \item Medir la latencia de ida y vuelta y el rendimiento con cargas de
    100 bytes y 10 KB (10240 bytes), mediante corridas repetidas y una prueba
    concurrente.
    \item Analizar la media, mediana, desviación estándar y percentil 95 de
    las latencias.
    \item Comprobar con registros reales que el reloj de recepción de cada
    mensaje sea mayor que su reloj de envío.
\end{enumerate}

\subsection{Hipótesis previa al experimento}

En condiciones equivalentes de ejecución en localhost, se espera que sockets
TCP con JSON presente una menor latencia media de ida y vuelta que gRPC con
Protocol Buffers para mensajes de 100 bytes, debido a su menor sobrecarga de
protocolo. Se espera que esta diferencia disminuya con mensajes de 10 KB
(10240 bytes), donde la serialización y transferencia de la carga pueden tener
mayor influencia.

Esta hipótesis es una predicción formulada antes de realizar las mediciones.
El contraste utiliza las latencias medias de ambos mecanismos para cada tamaño
de mensaje y modo de ejecución. Su interpretación se presenta en la discusión.
